\chapter*{ABSTRACT}
\addcontentsline{toc}{chapter}{Abstract}

This project presents a software system for real-time driver stress detection and adaptive
vehicle cabin environment control using physiological signals collected from a Wear OS
smartwatch. Building on the foundation of Patent Application No.\ ACN1408, which introduced
a wearable-to-vehicle framework for alcohol detection, the present work extends the platform
to continuous stress monitoring with automated environmental response.

The system acquires blood volume pulse (BVP), electrodermal activity (EDA), skin temperature,
accelerometer data, and inter-beat interval (IBI) from a wrist-worn device. A Bidirectional
Long Short-Term Memory network with an Attention mechanism classifies driver stress into four
levels---Calm, Mild, Moderate, and Critical---at 4~Hz with an inference latency below 50~ms.
The model was trained on the WESAD public wearable stress dataset and achieves an accuracy of
97.8\% on the three dominant stress classes with a weighted F1-score of 0.978. The quantized
TensorFlow Lite model occupies 72~KB, making it suitable for deployment on resource-constrained
wearable hardware.

Environmental responses are transmitted via a 12-byte Bluetooth Low Energy packet to an
Arduino Nano 33 BLE module, which controls RGB cabin lighting, ventilation fan speed, and
audio prompts according to the detected stress level. A simulation framework validates all
five system scenarios including a near-accident panic spike and gradual fatigue buildup.

\textbf{Keywords:} Driver Stress Detection, Machine Learning, Wearable Sensors, TensorFlow
Lite, Bluetooth Low Energy, Adaptive Vehicle Environment, BiLSTM, WESAD, Arduino
